% =====================================================================
%  MSACL DS301 Deep Learning · Quiz 14 · Agentic AI, Demo-Driven
%  Academic problem-set style (single-source, \ifsolution toggle)
%  SINGLE SOURCE — produces BOTH the student quiz and the answer key.
%
%  ANSWER TOGGLE
%  -------------
%  Student version (answers HIDDEN) — the default:
%       pdflatex quiz14_agents.tex
%  Answer key (answers SHOWN):
%       pdflatex "\def\showsolutions{}% =====================================================================
%  MSACL DS301 Deep Learning · Quiz 14 · Agentic AI, Demo-Driven
%  Academic problem-set style (single-source, \ifsolution toggle)
%  SINGLE SOURCE — produces BOTH the student quiz and the answer key.
%
%  ANSWER TOGGLE
%  -------------
%  Student version (answers HIDDEN) — the default:
%       pdflatex quiz14_agents.tex
%  Answer key (answers SHOWN):
%       pdflatex "\def\showsolutions{}% =====================================================================
%  MSACL DS301 Deep Learning · Quiz 14 · Agentic AI, Demo-Driven
%  Academic problem-set style (single-source, \ifsolution toggle)
%  SINGLE SOURCE — produces BOTH the student quiz and the answer key.
%
%  ANSWER TOGGLE
%  -------------
%  Student version (answers HIDDEN) — the default:
%       pdflatex quiz14_agents.tex
%  Answer key (answers SHOWN):
%       pdflatex "\def\showsolutions{}% =====================================================================
%  MSACL DS301 Deep Learning · Quiz 14 · Agentic AI, Demo-Driven
%  Academic problem-set style (single-source, \ifsolution toggle)
%  SINGLE SOURCE — produces BOTH the student quiz and the answer key.
%
%  ANSWER TOGGLE
%  -------------
%  Student version (answers HIDDEN) — the default:
%       pdflatex quiz14_agents.tex
%  Answer key (answers SHOWN):
%       pdflatex "\def\showsolutions{}\input{quiz14_agents.tex}"
%  (or, to flip by hand, change \solutionfalse to \solutiontrue below.)
% =====================================================================
\documentclass[11pt]{article}

\newif\ifsolution
\ifdefined\showsolutions \solutiontrue \else \solutionfalse \fi
% \solutiontrue   % <-- uncomment to hard-code the ANSWER KEY build

\usepackage[letterpaper, margin=0.6in, top=0.7in, bottom=0.5in]{geometry}
\usepackage{amsmath, amssymb}
\usepackage{xcolor}
\usepackage{array}
\usepackage{tikz}
\usetikzlibrary{arrows.meta, calc}
\usepackage{fancyhdr}

\pagestyle{fancy}
\fancyhf{}
\fancyhead[L]{\small\textbf{MSACL DS301 --- Deep Learning}}
\fancyhead[R]{\small\textbf{Quiz 14: Agentic AI\ifsolution\ \textmd{(Answer Key)}\fi}}
\fancyfoot[C]{\thepage}
\renewcommand{\headrulewidth}{0.4pt}

\newcommand{\blk}[1]{\underline{\hspace{#1}}}
% Reveal-or-blank: shows the answer in the key, a blank rule on the student sheet.
\newcommand{\rb}[2]{\ifsolution\textbf{#1}\else\blk{#2}\fi}
% Boxed answer (key only) and blank writing space (student only).
\newcommand{\keybox}[1]{\ifsolution\par\vspace{2pt}%
  \fbox{\parbox{0.965\textwidth}{\small\textbf{Answer.}\; #1}}\par\fi}
\newcommand{\work}[1]{\ifsolution\else\vspace{#1}\fi}

% ---- course palette ------------------------------------------------------
\definecolor{ink}{HTML}{1B2025}
\definecolor{teal}{HTML}{0E7C7B}
\definecolor{tealsoft}{HTML}{E3F0EF}
\definecolor{amber}{HTML}{E09E2F}
\definecolor{ambersoft}{HTML}{FBF0DC}
\definecolor{redd}{HTML}{C4534F}
\definecolor{redsoft}{HTML}{F7E4E3}
\definecolor{inksoft}{HTML}{3D444C}
\definecolor{muted}{HTML}{8A9099}

\setlength{\parindent}{0pt}
\setlength{\parskip}{3pt}
\setlength{\abovedisplayskip}{5pt}
\setlength{\belowdisplayskip}{5pt}

\begin{document}

\begin{center}
{\Large \textbf{Quiz 14 --- Agentic AI, Demo-Driven}}\\[2pt]
{\normalsize Label reason vs act (and spot a fabricated claim), match a guardrail to a new failure, and decide what stays human-reviewed}
\end{center}

\vspace{-2pt}
\begin{center}
\fbox{\parbox{0.92\textwidth}{\small
\textbf{Instructions:}\; Answer all three questions --- every part is answerable from the
ReAct pattern and the guardrails on the slides. No arithmetic; intuition only.%
\ifsolution\else\\[4pt]\textbf{Name:}\ \blk{6cm}\hfill\textbf{Date:}\ \blk{3.5cm}\fi
}}
\end{center}

\vspace{-2pt}\hrule\vspace{5pt}

% =========================== Q1 ===========================
\textbf{1. Label each step: reason or act?}\; {\small An agent runs the \textbf{ReAct} loop
(\textcolor{teal}{\textbf{Thought}} $\rightarrow$ \textcolor{amber}{\textbf{Action}} $\rightarrow$
Observation $\rightarrow \dots \rightarrow$ Answer). A step is \textbf{R} (\emph{reason}) if it is a
\textbf{Thought} --- the agent planning or interpreting --- and \textbf{A} (\emph{act}) if it is an
\textbf{Action} --- the agent calling a tool (read, compute, write). The grey observation lines are the
tool's replies, not steps. Below is one transcript (an assistant asked to summarise this week's
peptide-ID counts and note any underperforming run). \textbf{(a)} Write \textbf{R} or \textbf{A} beside
each numbered step.}

\vspace{5pt}
\begin{center}
\renewcommand{\arraystretch}{1.35}
\begin{tabular}{@{}r@{\ \ }p{0.72\textwidth}c@{}}
\textbf{1.} & ``Before I can compare runs, I need to load \texttt{results.csv} and see its columns.''
            & \rb{R}{1.4cm} \\
\textbf{2.} & \texttt{read\_csv("results.csv")} & \rb{A}{1.4cm} \\
\multicolumn{2}{@{}l}{\hspace{1.1em}\textcolor{muted}{\itshape $\rightarrow$ 6 runs $\cdot$ column: peptide\_ids}} & \\
\textbf{3.} & ``I'll compute the average ID count so I can see which runs fall below it.''
            & \rb{R}{1.4cm} \\
\textbf{4.} & compute the mean of \texttt{peptide\_ids} across the 6 runs & \rb{A}{1.4cm} \\
\multicolumn{2}{@{}l}{\hspace{1.1em}\textcolor{muted}{\itshape $\rightarrow$ mean $= 1{,}850$ $\cdot$ run B-04 $= 420$}} & \\
\textbf{5.} & ``B-04 (420) is far below the mean of 1,850, so that run underperformed.''
            & \rb{R}{1.4cm} \\
\textbf{6.} & \texttt{write\_summary(flag="B-04")} & \rb{A}{1.4cm} \\
\multicolumn{2}{@{}l}{\hspace{1.1em}\textcolor{muted}{\itshape $\rightarrow$ summary drafted}} & \\
\multicolumn{3}{@{}p{0.98\textwidth}}{\textcolor{redd}{\textbf{Answer:}}\ ``B-04 is low (420 vs mean
1{,}850). The cause is a \textbf{clogged emitter} --- I recommend replacing it.''} \\
\end{tabular}
\end{center}

\textbf{(b)} The final \textbf{Answer} states one thing that \emph{no observation ever showed}. Which
claim, and where should the agent have stopped for a human?

\work{0.9cm}
\ifsolution\else\blk{\linewidth}\fi

\keybox{\textbf{(a) 1 = R, 2 = A, 3 = R, 4 = A, 5 = R, 6 = A.}\; Odd steps are \emph{Thoughts} --- the
agent planning (1, 3) or interpreting a result (5) --- so they are \textbf{reason}. Even steps each call a
tool --- reading the file (2), computing the mean (4), writing the summary (6) --- so they are
\textbf{act}. The grey lines are \emph{observations} (what a tool returned), not agent steps.
\textbf{(b)} The claim ``\textbf{the cause is a clogged emitter}'' is \emph{unsupported} --- no tool ever
diagnosed a \emph{cause}; the observations show only that B-04's count is low, not \emph{why}. The agent
should have STOPPED after flagging the low run (step 6) and left any root-cause claim (and any repair) to a
human --- asserting an unverified cause is a hallucination that must not drive an action.}

\vspace{2pt}\hrule\vspace{5pt}

% =========================== Q2 ===========================
\textbf{2. Which guardrail for which failure?}\; {\small Each scenario below is a \textbf{new} failure ---
the wording won't match the slide line-for-line, so reason from what each guardrail \emph{does}. Match each
to the ONE guardrail that prevents it (each used once).}

\vspace{4pt}
\textbf{(a)} The agent's draft cites a \textbf{reference range that appears in no SOP}.
\hfill $\rightarrow$\;\rb{Human review}{4.6cm}

\vspace{3pt}
\textbf{(b)} Told to ``tidy old files,'' the agent \textbf{overwrites last month's raw \texttt{.raw} files}.
\hfill $\rightarrow$\;\rb{Restricted tools}{4.6cm}

\vspace{3pt}
\textbf{(c)} A released result was wrong and the lab \textbf{can't reconstruct which files the agent used}.
\hfill $\rightarrow$\;\rb{Logging}{4.6cm}

\vspace{4pt}
\begin{center}
{\small \textbf{Word bank:}\;\; Human review \;$\cdot$\; Restricted tools \;$\cdot$\; Logging}
\end{center}

\vspace{2pt}
\textbf{(d)} \emph{Human review} \textbf{catches} the fabricated citation in (a) --- but only if the
reviewer checks. Name \textbf{one} thing that would make such fabrications \textbf{rarer in the first
place}.

\work{0.9cm}
\ifsolution\else\blk{\linewidth}\fi

\keybox{(a) \textbf{Human review} --- a person verifies every number/citation against a source before
release. (b) \textbf{Restricted tools} --- read-only access, no write/delete/purchase without approval, so
it \emph{cannot} take a destructive action (the disguised version of ``deletes raw data''). (c)
\textbf{Logging} --- every reason$\rightarrow$act step recorded, so the run is auditable and reproducible.
Each used once. \textbf{(d)} Human review is a \emph{catch-after} control; to make fabrications rarer
\emph{up front}, \textbf{ground the agent in real retrieved sources (RAG)} and/or \textbf{restrict it to a
vetted reference database} so it can only cite documents that actually exist --- prevention layered on top
of detection.}

\vspace{2pt}\hrule\vspace{5pt}

% =========================== Q3 ===========================
\textbf{3. What must stay human-reviewed in a clinical workflow?}\; {\small An agent could help with each
step below. \textbf{Circle every step that a qualified human must review and sign before it reaches a
patient.}}

\vspace{4pt}
\begin{center}
\renewcommand{\arraystretch}{1.3}
\begin{tabular}{@{}ll@{}}
\textbf{(a)} \textbf{hold} a suspect run for a human to review &
\textbf{(b)} \textbf{auto-release} runs it judges ``in range'' \\
\textbf{(c)} file the final resistance (R/S) call in the report &
\textbf{(d)} reformat the QC table to the house template \\
\textbf{(e)} draft a plain-language note explaining a result & \\
\end{tabular}
\end{center}

\vspace{3pt}
Circle all that must stay human-reviewed:\; \rb{(b) and (c)}{5.5cm}

\vspace{3pt}
In one sentence, what is the rule of thumb?

\work{1.0cm}
\ifsolution\else\blk{\linewidth}\fi

\keybox{\textbf{(b) and (c)} must stay human-reviewed and signed. (b) auto-\emph{releasing} a result and
(c) the final resistance (R/S) call each become a \emph{released clinical result} or a
\emph{patient-affecting decision}, so a qualified person must verify and sign. Note the near-miss with (a):
\textbf{holding / flagging} a suspect run for review is \emph{conservative} --- it only routes work to a
human and releases nothing --- so an agent may do it freely, unlike auto-releasing in (b). Steps (d)
reformatting tables and (e) drafting an explanatory note are agent-friendly helpers (still spot-checked and
logged), because they don't release a result or decide anything about a patient. \textbf{Rule of thumb:}
anything that \emph{releases a result} or makes a \emph{patient-affecting decision} needs the human gate
--- flagging, drafting, and reformatting do not.}

\vspace{4pt}\hrule\vspace{3pt}
{\small\itshape \ifsolution Answer key.\else Keep this sheet --- next is the Lab 5 capstone: pick a track,
adapt a working pipeline, evaluate it properly, and write your first-DL-project one-pager.\fi}

\end{document}
"
%  (or, to flip by hand, change \solutionfalse to \solutiontrue below.)
% =====================================================================
\documentclass[11pt]{article}

\newif\ifsolution
\ifdefined\showsolutions \solutiontrue \else \solutionfalse \fi
% \solutiontrue   % <-- uncomment to hard-code the ANSWER KEY build

\usepackage[letterpaper, margin=0.6in, top=0.7in, bottom=0.5in]{geometry}
\usepackage{amsmath, amssymb}
\usepackage{xcolor}
\usepackage{array}
\usepackage{tikz}
\usetikzlibrary{arrows.meta, calc}
\usepackage{fancyhdr}

\pagestyle{fancy}
\fancyhf{}
\fancyhead[L]{\small\textbf{MSACL DS301 --- Deep Learning}}
\fancyhead[R]{\small\textbf{Quiz 14: Agentic AI\ifsolution\ \textmd{(Answer Key)}\fi}}
\fancyfoot[C]{\thepage}
\renewcommand{\headrulewidth}{0.4pt}

\newcommand{\blk}[1]{\underline{\hspace{#1}}}
% Reveal-or-blank: shows the answer in the key, a blank rule on the student sheet.
\newcommand{\rb}[2]{\ifsolution\textbf{#1}\else\blk{#2}\fi}
% Boxed answer (key only) and blank writing space (student only).
\newcommand{\keybox}[1]{\ifsolution\par\vspace{2pt}%
  \fbox{\parbox{0.965\textwidth}{\small\textbf{Answer.}\; #1}}\par\fi}
\newcommand{\work}[1]{\ifsolution\else\vspace{#1}\fi}

% ---- course palette ------------------------------------------------------
\definecolor{ink}{HTML}{1B2025}
\definecolor{teal}{HTML}{0E7C7B}
\definecolor{tealsoft}{HTML}{E3F0EF}
\definecolor{amber}{HTML}{E09E2F}
\definecolor{ambersoft}{HTML}{FBF0DC}
\definecolor{redd}{HTML}{C4534F}
\definecolor{redsoft}{HTML}{F7E4E3}
\definecolor{inksoft}{HTML}{3D444C}
\definecolor{muted}{HTML}{8A9099}

\setlength{\parindent}{0pt}
\setlength{\parskip}{3pt}
\setlength{\abovedisplayskip}{5pt}
\setlength{\belowdisplayskip}{5pt}

\begin{document}

\begin{center}
{\Large \textbf{Quiz 14 --- Agentic AI, Demo-Driven}}\\[2pt]
{\normalsize Label reason vs act (and spot a fabricated claim), match a guardrail to a new failure, and decide what stays human-reviewed}
\end{center}

\vspace{-2pt}
\begin{center}
\fbox{\parbox{0.92\textwidth}{\small
\textbf{Instructions:}\; Answer all three questions --- every part is answerable from the
ReAct pattern and the guardrails on the slides. No arithmetic; intuition only.%
\ifsolution\else\\[4pt]\textbf{Name:}\ \blk{6cm}\hfill\textbf{Date:}\ \blk{3.5cm}\fi
}}
\end{center}

\vspace{-2pt}\hrule\vspace{5pt}

% =========================== Q1 ===========================
\textbf{1. Label each step: reason or act?}\; {\small An agent runs the \textbf{ReAct} loop
(\textcolor{teal}{\textbf{Thought}} $\rightarrow$ \textcolor{amber}{\textbf{Action}} $\rightarrow$
Observation $\rightarrow \dots \rightarrow$ Answer). A step is \textbf{R} (\emph{reason}) if it is a
\textbf{Thought} --- the agent planning or interpreting --- and \textbf{A} (\emph{act}) if it is an
\textbf{Action} --- the agent calling a tool (read, compute, write). The grey observation lines are the
tool's replies, not steps. Below is one transcript (an assistant asked to summarise this week's
peptide-ID counts and note any underperforming run). \textbf{(a)} Write \textbf{R} or \textbf{A} beside
each numbered step.}

\vspace{5pt}
\begin{center}
\renewcommand{\arraystretch}{1.35}
\begin{tabular}{@{}r@{\ \ }p{0.72\textwidth}c@{}}
\textbf{1.} & ``Before I can compare runs, I need to load \texttt{results.csv} and see its columns.''
            & \rb{R}{1.4cm} \\
\textbf{2.} & \texttt{read\_csv("results.csv")} & \rb{A}{1.4cm} \\
\multicolumn{2}{@{}l}{\hspace{1.1em}\textcolor{muted}{\itshape $\rightarrow$ 6 runs $\cdot$ column: peptide\_ids}} & \\
\textbf{3.} & ``I'll compute the average ID count so I can see which runs fall below it.''
            & \rb{R}{1.4cm} \\
\textbf{4.} & compute the mean of \texttt{peptide\_ids} across the 6 runs & \rb{A}{1.4cm} \\
\multicolumn{2}{@{}l}{\hspace{1.1em}\textcolor{muted}{\itshape $\rightarrow$ mean $= 1{,}850$ $\cdot$ run B-04 $= 420$}} & \\
\textbf{5.} & ``B-04 (420) is far below the mean of 1,850, so that run underperformed.''
            & \rb{R}{1.4cm} \\
\textbf{6.} & \texttt{write\_summary(flag="B-04")} & \rb{A}{1.4cm} \\
\multicolumn{2}{@{}l}{\hspace{1.1em}\textcolor{muted}{\itshape $\rightarrow$ summary drafted}} & \\
\multicolumn{3}{@{}p{0.98\textwidth}}{\textcolor{redd}{\textbf{Answer:}}\ ``B-04 is low (420 vs mean
1{,}850). The cause is a \textbf{clogged emitter} --- I recommend replacing it.''} \\
\end{tabular}
\end{center}

\textbf{(b)} The final \textbf{Answer} states one thing that \emph{no observation ever showed}. Which
claim, and where should the agent have stopped for a human?

\work{0.9cm}
\ifsolution\else\blk{\linewidth}\fi

\keybox{\textbf{(a) 1 = R, 2 = A, 3 = R, 4 = A, 5 = R, 6 = A.}\; Odd steps are \emph{Thoughts} --- the
agent planning (1, 3) or interpreting a result (5) --- so they are \textbf{reason}. Even steps each call a
tool --- reading the file (2), computing the mean (4), writing the summary (6) --- so they are
\textbf{act}. The grey lines are \emph{observations} (what a tool returned), not agent steps.
\textbf{(b)} The claim ``\textbf{the cause is a clogged emitter}'' is \emph{unsupported} --- no tool ever
diagnosed a \emph{cause}; the observations show only that B-04's count is low, not \emph{why}. The agent
should have STOPPED after flagging the low run (step 6) and left any root-cause claim (and any repair) to a
human --- asserting an unverified cause is a hallucination that must not drive an action.}

\vspace{2pt}\hrule\vspace{5pt}

% =========================== Q2 ===========================
\textbf{2. Which guardrail for which failure?}\; {\small Each scenario below is a \textbf{new} failure ---
the wording won't match the slide line-for-line, so reason from what each guardrail \emph{does}. Match each
to the ONE guardrail that prevents it (each used once).}

\vspace{4pt}
\textbf{(a)} The agent's draft cites a \textbf{reference range that appears in no SOP}.
\hfill $\rightarrow$\;\rb{Human review}{4.6cm}

\vspace{3pt}
\textbf{(b)} Told to ``tidy old files,'' the agent \textbf{overwrites last month's raw \texttt{.raw} files}.
\hfill $\rightarrow$\;\rb{Restricted tools}{4.6cm}

\vspace{3pt}
\textbf{(c)} A released result was wrong and the lab \textbf{can't reconstruct which files the agent used}.
\hfill $\rightarrow$\;\rb{Logging}{4.6cm}

\vspace{4pt}
\begin{center}
{\small \textbf{Word bank:}\;\; Human review \;$\cdot$\; Restricted tools \;$\cdot$\; Logging}
\end{center}

\vspace{2pt}
\textbf{(d)} \emph{Human review} \textbf{catches} the fabricated citation in (a) --- but only if the
reviewer checks. Name \textbf{one} thing that would make such fabrications \textbf{rarer in the first
place}.

\work{0.9cm}
\ifsolution\else\blk{\linewidth}\fi

\keybox{(a) \textbf{Human review} --- a person verifies every number/citation against a source before
release. (b) \textbf{Restricted tools} --- read-only access, no write/delete/purchase without approval, so
it \emph{cannot} take a destructive action (the disguised version of ``deletes raw data''). (c)
\textbf{Logging} --- every reason$\rightarrow$act step recorded, so the run is auditable and reproducible.
Each used once. \textbf{(d)} Human review is a \emph{catch-after} control; to make fabrications rarer
\emph{up front}, \textbf{ground the agent in real retrieved sources (RAG)} and/or \textbf{restrict it to a
vetted reference database} so it can only cite documents that actually exist --- prevention layered on top
of detection.}

\vspace{2pt}\hrule\vspace{5pt}

% =========================== Q3 ===========================
\textbf{3. What must stay human-reviewed in a clinical workflow?}\; {\small An agent could help with each
step below. \textbf{Circle every step that a qualified human must review and sign before it reaches a
patient.}}

\vspace{4pt}
\begin{center}
\renewcommand{\arraystretch}{1.3}
\begin{tabular}{@{}ll@{}}
\textbf{(a)} \textbf{hold} a suspect run for a human to review &
\textbf{(b)} \textbf{auto-release} runs it judges ``in range'' \\
\textbf{(c)} file the final resistance (R/S) call in the report &
\textbf{(d)} reformat the QC table to the house template \\
\textbf{(e)} draft a plain-language note explaining a result & \\
\end{tabular}
\end{center}

\vspace{3pt}
Circle all that must stay human-reviewed:\; \rb{(b) and (c)}{5.5cm}

\vspace{3pt}
In one sentence, what is the rule of thumb?

\work{1.0cm}
\ifsolution\else\blk{\linewidth}\fi

\keybox{\textbf{(b) and (c)} must stay human-reviewed and signed. (b) auto-\emph{releasing} a result and
(c) the final resistance (R/S) call each become a \emph{released clinical result} or a
\emph{patient-affecting decision}, so a qualified person must verify and sign. Note the near-miss with (a):
\textbf{holding / flagging} a suspect run for review is \emph{conservative} --- it only routes work to a
human and releases nothing --- so an agent may do it freely, unlike auto-releasing in (b). Steps (d)
reformatting tables and (e) drafting an explanatory note are agent-friendly helpers (still spot-checked and
logged), because they don't release a result or decide anything about a patient. \textbf{Rule of thumb:}
anything that \emph{releases a result} or makes a \emph{patient-affecting decision} needs the human gate
--- flagging, drafting, and reformatting do not.}

\vspace{4pt}\hrule\vspace{3pt}
{\small\itshape \ifsolution Answer key.\else Keep this sheet --- next is the Lab 5 capstone: pick a track,
adapt a working pipeline, evaluate it properly, and write your first-DL-project one-pager.\fi}

\end{document}
"
%  (or, to flip by hand, change \solutionfalse to \solutiontrue below.)
% =====================================================================
\documentclass[11pt]{article}

\newif\ifsolution
\ifdefined\showsolutions \solutiontrue \else \solutionfalse \fi
% \solutiontrue   % <-- uncomment to hard-code the ANSWER KEY build

\usepackage[letterpaper, margin=0.6in, top=0.7in, bottom=0.5in]{geometry}
\usepackage{amsmath, amssymb}
\usepackage{xcolor}
\usepackage{array}
\usepackage{tikz}
\usetikzlibrary{arrows.meta, calc}
\usepackage{fancyhdr}

\pagestyle{fancy}
\fancyhf{}
\fancyhead[L]{\small\textbf{MSACL DS301 --- Deep Learning}}
\fancyhead[R]{\small\textbf{Quiz 14: Agentic AI\ifsolution\ \textmd{(Answer Key)}\fi}}
\fancyfoot[C]{\thepage}
\renewcommand{\headrulewidth}{0.4pt}

\newcommand{\blk}[1]{\underline{\hspace{#1}}}
% Reveal-or-blank: shows the answer in the key, a blank rule on the student sheet.
\newcommand{\rb}[2]{\ifsolution\textbf{#1}\else\blk{#2}\fi}
% Boxed answer (key only) and blank writing space (student only).
\newcommand{\keybox}[1]{\ifsolution\par\vspace{2pt}%
  \fbox{\parbox{0.965\textwidth}{\small\textbf{Answer.}\; #1}}\par\fi}
\newcommand{\work}[1]{\ifsolution\else\vspace{#1}\fi}

% ---- course palette ------------------------------------------------------
\definecolor{ink}{HTML}{1B2025}
\definecolor{teal}{HTML}{0E7C7B}
\definecolor{tealsoft}{HTML}{E3F0EF}
\definecolor{amber}{HTML}{E09E2F}
\definecolor{ambersoft}{HTML}{FBF0DC}
\definecolor{redd}{HTML}{C4534F}
\definecolor{redsoft}{HTML}{F7E4E3}
\definecolor{inksoft}{HTML}{3D444C}
\definecolor{muted}{HTML}{8A9099}

\setlength{\parindent}{0pt}
\setlength{\parskip}{3pt}
\setlength{\abovedisplayskip}{5pt}
\setlength{\belowdisplayskip}{5pt}

\begin{document}

\begin{center}
{\Large \textbf{Quiz 14 --- Agentic AI, Demo-Driven}}\\[2pt]
{\normalsize Label reason vs act (and spot a fabricated claim), match a guardrail to a new failure, and decide what stays human-reviewed}
\end{center}

\vspace{-2pt}
\begin{center}
\fbox{\parbox{0.92\textwidth}{\small
\textbf{Instructions:}\; Answer all three questions --- every part is answerable from the
ReAct pattern and the guardrails on the slides. No arithmetic; intuition only.%
\ifsolution\else\\[4pt]\textbf{Name:}\ \blk{6cm}\hfill\textbf{Date:}\ \blk{3.5cm}\fi
}}
\end{center}

\vspace{-2pt}\hrule\vspace{5pt}

% =========================== Q1 ===========================
\textbf{1. Label each step: reason or act?}\; {\small An agent runs the \textbf{ReAct} loop
(\textcolor{teal}{\textbf{Thought}} $\rightarrow$ \textcolor{amber}{\textbf{Action}} $\rightarrow$
Observation $\rightarrow \dots \rightarrow$ Answer). A step is \textbf{R} (\emph{reason}) if it is a
\textbf{Thought} --- the agent planning or interpreting --- and \textbf{A} (\emph{act}) if it is an
\textbf{Action} --- the agent calling a tool (read, compute, write). The grey observation lines are the
tool's replies, not steps. Below is one transcript (an assistant asked to summarise this week's
peptide-ID counts and note any underperforming run). \textbf{(a)} Write \textbf{R} or \textbf{A} beside
each numbered step.}

\vspace{5pt}
\begin{center}
\renewcommand{\arraystretch}{1.35}
\begin{tabular}{@{}r@{\ \ }p{0.72\textwidth}c@{}}
\textbf{1.} & ``Before I can compare runs, I need to load \texttt{results.csv} and see its columns.''
            & \rb{R}{1.4cm} \\
\textbf{2.} & \texttt{read\_csv("results.csv")} & \rb{A}{1.4cm} \\
\multicolumn{2}{@{}l}{\hspace{1.1em}\textcolor{muted}{\itshape $\rightarrow$ 6 runs $\cdot$ column: peptide\_ids}} & \\
\textbf{3.} & ``I'll compute the average ID count so I can see which runs fall below it.''
            & \rb{R}{1.4cm} \\
\textbf{4.} & compute the mean of \texttt{peptide\_ids} across the 6 runs & \rb{A}{1.4cm} \\
\multicolumn{2}{@{}l}{\hspace{1.1em}\textcolor{muted}{\itshape $\rightarrow$ mean $= 1{,}850$ $\cdot$ run B-04 $= 420$}} & \\
\textbf{5.} & ``B-04 (420) is far below the mean of 1,850, so that run underperformed.''
            & \rb{R}{1.4cm} \\
\textbf{6.} & \texttt{write\_summary(flag="B-04")} & \rb{A}{1.4cm} \\
\multicolumn{2}{@{}l}{\hspace{1.1em}\textcolor{muted}{\itshape $\rightarrow$ summary drafted}} & \\
\multicolumn{3}{@{}p{0.98\textwidth}}{\textcolor{redd}{\textbf{Answer:}}\ ``B-04 is low (420 vs mean
1{,}850). The cause is a \textbf{clogged emitter} --- I recommend replacing it.''} \\
\end{tabular}
\end{center}

\textbf{(b)} The final \textbf{Answer} states one thing that \emph{no observation ever showed}. Which
claim, and where should the agent have stopped for a human?

\work{0.9cm}
\ifsolution\else\blk{\linewidth}\fi

\keybox{\textbf{(a) 1 = R, 2 = A, 3 = R, 4 = A, 5 = R, 6 = A.}\; Odd steps are \emph{Thoughts} --- the
agent planning (1, 3) or interpreting a result (5) --- so they are \textbf{reason}. Even steps each call a
tool --- reading the file (2), computing the mean (4), writing the summary (6) --- so they are
\textbf{act}. The grey lines are \emph{observations} (what a tool returned), not agent steps.
\textbf{(b)} The claim ``\textbf{the cause is a clogged emitter}'' is \emph{unsupported} --- no tool ever
diagnosed a \emph{cause}; the observations show only that B-04's count is low, not \emph{why}. The agent
should have STOPPED after flagging the low run (step 6) and left any root-cause claim (and any repair) to a
human --- asserting an unverified cause is a hallucination that must not drive an action.}

\vspace{2pt}\hrule\vspace{5pt}

% =========================== Q2 ===========================
\textbf{2. Which guardrail for which failure?}\; {\small Each scenario below is a \textbf{new} failure ---
the wording won't match the slide line-for-line, so reason from what each guardrail \emph{does}. Match each
to the ONE guardrail that prevents it (each used once).}

\vspace{4pt}
\textbf{(a)} The agent's draft cites a \textbf{reference range that appears in no SOP}.
\hfill $\rightarrow$\;\rb{Human review}{4.6cm}

\vspace{3pt}
\textbf{(b)} Told to ``tidy old files,'' the agent \textbf{overwrites last month's raw \texttt{.raw} files}.
\hfill $\rightarrow$\;\rb{Restricted tools}{4.6cm}

\vspace{3pt}
\textbf{(c)} A released result was wrong and the lab \textbf{can't reconstruct which files the agent used}.
\hfill $\rightarrow$\;\rb{Logging}{4.6cm}

\vspace{4pt}
\begin{center}
{\small \textbf{Word bank:}\;\; Human review \;$\cdot$\; Restricted tools \;$\cdot$\; Logging}
\end{center}

\vspace{2pt}
\textbf{(d)} \emph{Human review} \textbf{catches} the fabricated citation in (a) --- but only if the
reviewer checks. Name \textbf{one} thing that would make such fabrications \textbf{rarer in the first
place}.

\work{0.9cm}
\ifsolution\else\blk{\linewidth}\fi

\keybox{(a) \textbf{Human review} --- a person verifies every number/citation against a source before
release. (b) \textbf{Restricted tools} --- read-only access, no write/delete/purchase without approval, so
it \emph{cannot} take a destructive action (the disguised version of ``deletes raw data''). (c)
\textbf{Logging} --- every reason$\rightarrow$act step recorded, so the run is auditable and reproducible.
Each used once. \textbf{(d)} Human review is a \emph{catch-after} control; to make fabrications rarer
\emph{up front}, \textbf{ground the agent in real retrieved sources (RAG)} and/or \textbf{restrict it to a
vetted reference database} so it can only cite documents that actually exist --- prevention layered on top
of detection.}

\vspace{2pt}\hrule\vspace{5pt}

% =========================== Q3 ===========================
\textbf{3. What must stay human-reviewed in a clinical workflow?}\; {\small An agent could help with each
step below. \textbf{Circle every step that a qualified human must review and sign before it reaches a
patient.}}

\vspace{4pt}
\begin{center}
\renewcommand{\arraystretch}{1.3}
\begin{tabular}{@{}ll@{}}
\textbf{(a)} \textbf{hold} a suspect run for a human to review &
\textbf{(b)} \textbf{auto-release} runs it judges ``in range'' \\
\textbf{(c)} file the final resistance (R/S) call in the report &
\textbf{(d)} reformat the QC table to the house template \\
\textbf{(e)} draft a plain-language note explaining a result & \\
\end{tabular}
\end{center}

\vspace{3pt}
Circle all that must stay human-reviewed:\; \rb{(b) and (c)}{5.5cm}

\vspace{3pt}
In one sentence, what is the rule of thumb?

\work{1.0cm}
\ifsolution\else\blk{\linewidth}\fi

\keybox{\textbf{(b) and (c)} must stay human-reviewed and signed. (b) auto-\emph{releasing} a result and
(c) the final resistance (R/S) call each become a \emph{released clinical result} or a
\emph{patient-affecting decision}, so a qualified person must verify and sign. Note the near-miss with (a):
\textbf{holding / flagging} a suspect run for review is \emph{conservative} --- it only routes work to a
human and releases nothing --- so an agent may do it freely, unlike auto-releasing in (b). Steps (d)
reformatting tables and (e) drafting an explanatory note are agent-friendly helpers (still spot-checked and
logged), because they don't release a result or decide anything about a patient. \textbf{Rule of thumb:}
anything that \emph{releases a result} or makes a \emph{patient-affecting decision} needs the human gate
--- flagging, drafting, and reformatting do not.}

\vspace{4pt}\hrule\vspace{3pt}
{\small\itshape \ifsolution Answer key.\else Keep this sheet --- next is the Lab 5 capstone: pick a track,
adapt a working pipeline, evaluate it properly, and write your first-DL-project one-pager.\fi}

\end{document}
"
%  (or, to flip by hand, change \solutionfalse to \solutiontrue below.)
% =====================================================================
\documentclass[11pt]{article}

\newif\ifsolution
\ifdefined\showsolutions \solutiontrue \else \solutionfalse \fi
% \solutiontrue   % <-- uncomment to hard-code the ANSWER KEY build

\usepackage[letterpaper, margin=0.6in, top=0.7in, bottom=0.5in]{geometry}
\usepackage{amsmath, amssymb}
\usepackage{xcolor}
\usepackage{array}
\usepackage{tikz}
\usetikzlibrary{arrows.meta, calc}
\usepackage{fancyhdr}

\pagestyle{fancy}
\fancyhf{}
\fancyhead[L]{\small\textbf{MSACL DS301 --- Deep Learning}}
\fancyhead[R]{\small\textbf{Quiz 14: Agentic AI\ifsolution\ \textmd{(Answer Key)}\fi}}
\fancyfoot[C]{\thepage}
\renewcommand{\headrulewidth}{0.4pt}

\newcommand{\blk}[1]{\underline{\hspace{#1}}}
% Reveal-or-blank: shows the answer in the key, a blank rule on the student sheet.
\newcommand{\rb}[2]{\ifsolution\textbf{#1}\else\blk{#2}\fi}
% Boxed answer (key only) and blank writing space (student only).
\newcommand{\keybox}[1]{\ifsolution\par\vspace{2pt}%
  \fbox{\parbox{0.965\textwidth}{\small\textbf{Answer.}\; #1}}\par\fi}
\newcommand{\work}[1]{\ifsolution\else\vspace{#1}\fi}

% ---- course palette ------------------------------------------------------
\definecolor{ink}{HTML}{1B2025}
\definecolor{teal}{HTML}{0E7C7B}
\definecolor{tealsoft}{HTML}{E3F0EF}
\definecolor{amber}{HTML}{E09E2F}
\definecolor{ambersoft}{HTML}{FBF0DC}
\definecolor{redd}{HTML}{C4534F}
\definecolor{redsoft}{HTML}{F7E4E3}
\definecolor{inksoft}{HTML}{3D444C}
\definecolor{muted}{HTML}{8A9099}

\setlength{\parindent}{0pt}
\setlength{\parskip}{3pt}
\setlength{\abovedisplayskip}{5pt}
\setlength{\belowdisplayskip}{5pt}

\begin{document}

\begin{center}
{\Large \textbf{Quiz 14 --- Agentic AI, Demo-Driven}}\\[2pt]
{\normalsize Label reason vs act (and spot a fabricated claim), match a guardrail to a new failure, and decide what stays human-reviewed}
\end{center}

\vspace{-2pt}
\begin{center}
\fbox{\parbox{0.92\textwidth}{\small
\textbf{Instructions:}\; Answer all three questions --- every part is answerable from the
ReAct pattern and the guardrails on the slides. No arithmetic; intuition only.%
\ifsolution\else\\[4pt]\textbf{Name:}\ \blk{6cm}\hfill\textbf{Date:}\ \blk{3.5cm}\fi
}}
\end{center}

\vspace{-2pt}\hrule\vspace{5pt}

% =========================== Q1 ===========================
\textbf{1. Label each step: reason or act?}\; {\small An agent runs the \textbf{ReAct} loop
(\textcolor{teal}{\textbf{Thought}} $\rightarrow$ \textcolor{amber}{\textbf{Action}} $\rightarrow$
Observation $\rightarrow \dots \rightarrow$ Answer). A step is \textbf{R} (\emph{reason}) if it is a
\textbf{Thought} --- the agent planning or interpreting --- and \textbf{A} (\emph{act}) if it is an
\textbf{Action} --- the agent calling a tool (read, compute, write). The grey observation lines are the
tool's replies, not steps. Below is one transcript (an assistant asked to summarise this week's
peptide-ID counts and note any underperforming run). \textbf{(a)} Write \textbf{R} or \textbf{A} beside
each numbered step.}

\vspace{5pt}
\begin{center}
\renewcommand{\arraystretch}{1.35}
\begin{tabular}{@{}r@{\ \ }p{0.72\textwidth}c@{}}
\textbf{1.} & ``Before I can compare runs, I need to load \texttt{results.csv} and see its columns.''
            & \rb{R}{1.4cm} \\
\textbf{2.} & \texttt{read\_csv("results.csv")} & \rb{A}{1.4cm} \\
\multicolumn{2}{@{}l}{\hspace{1.1em}\textcolor{muted}{\itshape $\rightarrow$ 6 runs $\cdot$ column: peptide\_ids}} & \\
\textbf{3.} & ``I'll compute the average ID count so I can see which runs fall below it.''
            & \rb{R}{1.4cm} \\
\textbf{4.} & compute the mean of \texttt{peptide\_ids} across the 6 runs & \rb{A}{1.4cm} \\
\multicolumn{2}{@{}l}{\hspace{1.1em}\textcolor{muted}{\itshape $\rightarrow$ mean $= 1{,}850$ $\cdot$ run B-04 $= 420$}} & \\
\textbf{5.} & ``B-04 (420) is far below the mean of 1,850, so that run underperformed.''
            & \rb{R}{1.4cm} \\
\textbf{6.} & \texttt{write\_summary(flag="B-04")} & \rb{A}{1.4cm} \\
\multicolumn{2}{@{}l}{\hspace{1.1em}\textcolor{muted}{\itshape $\rightarrow$ summary drafted}} & \\
\multicolumn{3}{@{}p{0.98\textwidth}}{\textcolor{redd}{\textbf{Answer:}}\ ``B-04 is low (420 vs mean
1{,}850). The cause is a \textbf{clogged emitter} --- I recommend replacing it.''} \\
\end{tabular}
\end{center}

\textbf{(b)} The final \textbf{Answer} states one thing that \emph{no observation ever showed}. Which
claim, and where should the agent have stopped for a human?

\work{0.9cm}
\ifsolution\else\blk{\linewidth}\fi

\keybox{\textbf{(a) 1 = R, 2 = A, 3 = R, 4 = A, 5 = R, 6 = A.}\; Odd steps are \emph{Thoughts} --- the
agent planning (1, 3) or interpreting a result (5) --- so they are \textbf{reason}. Even steps each call a
tool --- reading the file (2), computing the mean (4), writing the summary (6) --- so they are
\textbf{act}. The grey lines are \emph{observations} (what a tool returned), not agent steps.
\textbf{(b)} The claim ``\textbf{the cause is a clogged emitter}'' is \emph{unsupported} --- no tool ever
diagnosed a \emph{cause}; the observations show only that B-04's count is low, not \emph{why}. The agent
should have STOPPED after flagging the low run (step 6) and left any root-cause claim (and any repair) to a
human --- asserting an unverified cause is a hallucination that must not drive an action.}

\vspace{2pt}\hrule\vspace{5pt}

% =========================== Q2 ===========================
\textbf{2. Which guardrail for which failure?}\; {\small Each scenario below is a \textbf{new} failure ---
the wording won't match the slide line-for-line, so reason from what each guardrail \emph{does}. Match each
to the ONE guardrail that prevents it (each used once).}

\vspace{4pt}
\textbf{(a)} The agent's draft cites a \textbf{reference range that appears in no SOP}.
\hfill $\rightarrow$\;\rb{Human review}{4.6cm}

\vspace{3pt}
\textbf{(b)} Told to ``tidy old files,'' the agent \textbf{overwrites last month's raw \texttt{.raw} files}.
\hfill $\rightarrow$\;\rb{Restricted tools}{4.6cm}

\vspace{3pt}
\textbf{(c)} A released result was wrong and the lab \textbf{can't reconstruct which files the agent used}.
\hfill $\rightarrow$\;\rb{Logging}{4.6cm}

\vspace{4pt}
\begin{center}
{\small \textbf{Word bank:}\;\; Human review \;$\cdot$\; Restricted tools \;$\cdot$\; Logging}
\end{center}

\vspace{2pt}
\textbf{(d)} \emph{Human review} \textbf{catches} the fabricated citation in (a) --- but only if the
reviewer checks. Name \textbf{one} thing that would make such fabrications \textbf{rarer in the first
place}.

\work{0.9cm}
\ifsolution\else\blk{\linewidth}\fi

\keybox{(a) \textbf{Human review} --- a person verifies every number/citation against a source before
release. (b) \textbf{Restricted tools} --- read-only access, no write/delete/purchase without approval, so
it \emph{cannot} take a destructive action (the disguised version of ``deletes raw data''). (c)
\textbf{Logging} --- every reason$\rightarrow$act step recorded, so the run is auditable and reproducible.
Each used once. \textbf{(d)} Human review is a \emph{catch-after} control; to make fabrications rarer
\emph{up front}, \textbf{ground the agent in real retrieved sources (RAG)} and/or \textbf{restrict it to a
vetted reference database} so it can only cite documents that actually exist --- prevention layered on top
of detection.}

\vspace{2pt}\hrule\vspace{5pt}

% =========================== Q3 ===========================
\textbf{3. What must stay human-reviewed in a clinical workflow?}\; {\small An agent could help with each
step below. \textbf{Circle every step that a qualified human must review and sign before it reaches a
patient.}}

\vspace{4pt}
\begin{center}
\renewcommand{\arraystretch}{1.3}
\begin{tabular}{@{}ll@{}}
\textbf{(a)} \textbf{hold} a suspect run for a human to review &
\textbf{(b)} \textbf{auto-release} runs it judges ``in range'' \\
\textbf{(c)} file the final resistance (R/S) call in the report &
\textbf{(d)} reformat the QC table to the house template \\
\textbf{(e)} draft a plain-language note explaining a result & \\
\end{tabular}
\end{center}

\vspace{3pt}
Circle all that must stay human-reviewed:\; \rb{(b) and (c)}{5.5cm}

\vspace{3pt}
In one sentence, what is the rule of thumb?

\work{1.0cm}
\ifsolution\else\blk{\linewidth}\fi

\keybox{\textbf{(b) and (c)} must stay human-reviewed and signed. (b) auto-\emph{releasing} a result and
(c) the final resistance (R/S) call each become a \emph{released clinical result} or a
\emph{patient-affecting decision}, so a qualified person must verify and sign. Note the near-miss with (a):
\textbf{holding / flagging} a suspect run for review is \emph{conservative} --- it only routes work to a
human and releases nothing --- so an agent may do it freely, unlike auto-releasing in (b). Steps (d)
reformatting tables and (e) drafting an explanatory note are agent-friendly helpers (still spot-checked and
logged), because they don't release a result or decide anything about a patient. \textbf{Rule of thumb:}
anything that \emph{releases a result} or makes a \emph{patient-affecting decision} needs the human gate
--- flagging, drafting, and reformatting do not.}

\vspace{4pt}\hrule\vspace{3pt}
{\small\itshape \ifsolution Answer key.\else Keep this sheet --- next is the Lab 5 capstone: pick a track,
adapt a working pipeline, evaluate it properly, and write your first-DL-project one-pager.\fi}

\end{document}
